\documentclass[11pt]{article}
\usepackage[margin=1in]{geometry}
\usepackage{amsmath,amssymb,amsthm,hyperref,enumitem}
\newtheorem{theorem}{Theorem}
\newtheorem{lemma}{Lemma}
\newcommand{\e}{\mathbf e}
\title{Positivity and monotonicity of the Dilcher--Ulas sequence}
\author{Principia Math In Progress}
\date{18 September 2026}
\begin{document}
\maketitle
\begin{abstract}
We prove the positivity and monotonicity conjecture for
$f_{2^k,n}(-1)$ in every interval between its prescribed zeros.
A finite-path representation makes each signed term, and each required
successive difference, a positive sum of numbers of walks. The argument
also gives the equality of the final two terms of every block.
\end{abstract}

\section*{Task 1: complete proof}
Dilcher and Ulas define
\[
 f_{m,n}(z)=\sum_{a=0}^n\binom na z^{\binom am},
 \qquad f_{m,n}=f_{m,n}(-1).
\]
The target is Conjecture 5.8 in the current preprint \cite{DUa},
and Conjecture 2 in the published paper \cite{DU}.

\begin{theorem}\label{main}
Let $k\ge1$, $m=2^k$, and $\nu\ge0$ be integers. Then
\[
 0<(-1)^\nu f_{m,2m\nu}<(-1)^\nu f_{m,2m\nu+1}
 <\cdots<(-1)^\nu f_{m,2m\nu+2m-3}
 =(-1)^\nu f_{m,2m\nu+2m-2}.
\]
Moreover, $f_{m,2m\nu+2m-1}=0$. Thus the displayed blocks
contain all nonzero terms of the sequence.
\end{theorem}

\subsection*{From the binomial sum to a path}
Lucas's parity rule gives
$\binom a{2^k}\equiv\lfloor a/2^k\rfloor\pmod2$.
For completeness, in $\mathbb F_2[x]$ write
$(1+x)^a=\prod_{r:a_r=1}(1+x^{2^r})$, where $a_r$ are the
binary digits of $a$. The coefficient of $x^{2^k}$ is $a_k$:
the smaller binary powers sum to less than $2^k$.
Consequently it suffices to study
\begin{equation}\label{F}
 F_m(n)=\sum_{a=0}^n\binom na(-1)^{\lfloor a/m\rfloor}.
\end{equation}
In fact, we shall prove the same block assertion for $F_m$ for
\emph{every integer $m\ge2$}. For powers of two, $F_m(n)=f_{m,n}$.

Let $B$ be the adjacency matrix of the path on vertices
$1,2,\ldots,2m-1$, let $\e=\e_m$ be the central unit vector, and define
\[
 U_0(x)=1,\quad U_1(x)=2x,\quad
 U_{t+1}(x)=2xU_t(x)-U_{t-1}(x),\qquad
 u_t=U_t(B/2)\e.
\]
Thus $u_0=\e$, $u_1=B\e$, and $u_{t+1}=Bu_t-u_{t-1}$.

\begin{lemma}\label{spectral}
For every $n,t\ge0$,
\begin{equation}\label{matrix}
 F_m(n)=\e^{\mathsf T}B^n u_n,
 \qquad u_{t+2m}=-u_t.
\end{equation}
\end{lemma}
\begin{proof}
Put $\theta_\ell=(2\ell+1)\pi/(2m)$ and
$z_\ell=e^{2i\theta_\ell}$, $0\le\ell<m$.
The sequence $\sigma(a)=(-1)^{\lfloor a/m\rfloor}$ is
antiperiodic with shift $m$. Finite Fourier inversion on its first
$m$ terms gives
\[
 \sigma(a)=\sum_{\ell=0}^{m-1}c_\ell z_\ell^a,
 \qquad c_\ell=\frac1m\sum_{b=0}^{m-1}z_\ell^{-b}
 =\frac{2}{m(1-z_\ell^{-1})}
 =\frac{-i e^{i\theta_\ell}}{m\sin\theta_\ell}.
\]
Indeed, the $z_\ell^a$ are orthogonal on $0\le a<m$ and all
have the same antiperiodicity. Apply the binomial theorem and take
real parts, using $1+z_\ell=2\cos\theta_\ell\,e^{i\theta_\ell}$:
\begin{equation}\label{trig}
 F_m(n)=\frac1m\sum_{\ell=0}^{m-1}
 (2\cos\theta_\ell)^n
 \frac{\sin((n+1)\theta_\ell)}{\sin\theta_\ell}.
\end{equation}
For powers of two, pairing the terms with opposite cosines recovers
Dilcher--Ulas Proposition 7, equation (38), or preprint
Proposition 5.3, equation (5.6). That trigonometric identity is known;
the derivation above fixes the normalization and also handles general $m$.

The normalized eigenvectors of $B$ have coordinates
\[
 v_r(a)=m^{-1/2}\sin\frac{ar\pi}{2m},\quad
 1\le r,a\le2m-1,
 \qquad Bv_r=2\cos\frac{r\pi}{2m}\,v_r.
\]
The sine recurrence verifies the eigenvalue equation, including the
zero boundary values. Distinct eigenvalues give orthogonality, and
$\sum_{a=1}^{2m-1}\sin^2(ar\pi/(2m))=m$ gives the normalization.
At the central vertex $v_r(m)^2$ is $1/m$ for odd $r$ and zero
for even $r$. Since $U_t(\cos\theta)=\sin((t+1)\theta)/\sin\theta$,
the spectral expansion of $\e^{\mathsf T}B^nU_n(B/2)\e$
is precisely \eqref{trig}.
On every odd mode, shifting $t$ by $2m$ changes the sine's sign.
The even modes have zero coefficient in $\e$, proving the second
identity in \eqref{matrix} as a vector identity.
\end{proof}

\subsection*{An explicit nonnegative central column}
For $0\le s\le m-1$, set
\[
 S_s=\sum_{a=0}^s\e_{m-s+2a},\qquad S_{-1}=0.
\]
\begin{lemma}\label{wave}
For $0\le t\le2m-2$,
\begin{equation}\label{column}
 u_t=S_{\min(t,\,2m-2-t)},\qquad u_{2m-1}=0.
\end{equation}
In particular, $u_t$ is nonnegative. If $t$ is even it contains
$\e_m$ with coefficient one; if $t$ is odd it contains
$\e_{m-1}$ with coefficient one.
\end{lemma}
\begin{proof}
Moving one step left or right on the path gives
\[
 BS_s=S_{s-1}+S_{s+1}\quad(0\le s\le m-2),
 \qquad BS_{m-1}=2S_{m-2}.
\]
These identities can also be checked at each coordinate: the two
outer coordinates have multiplicity one until the endpoints are
reached, and all overlapping coordinates have multiplicity two.
Starting with $u_0=S_0$ and $u_1=S_1$, the recurrence
$u_{t+1}=Bu_t-u_{t-1}$ first gives
$S_0,S_1,\ldots,S_{m-1}$ and then
$S_{m-2},\ldots,S_0,0$. This proves \eqref{column}, including $m=2$.
The minimum in \eqref{column} has the parity of $t$.
Its defining sum therefore contains the stated central vertex or
neighbor. In the odd case this minimum is at least one.
\end{proof}

\begin{lemma}\label{positive}
If $0\le t\le2m-2$, $n\ge0$, and $n\equiv t\pmod2$, then
$\e^{\mathsf T}B^n u_t>0$.
\end{lemma}
\begin{proof}
Every entry of $B^n$ counts walks of length $n$ on the path and
is nonnegative. When $n$ is even, there is a walk from $m$ to
$m$: for $n=0$ take the empty walk, and otherwise repeat the
edge between $m$ and $m-1$ back and forth. When $n$ is odd,
there is similarly a walk from $m$ to $m-1$. The corresponding
coordinate of $u_t$ equals one by Lemma~\ref{wave}.
\end{proof}

\subsection*{All blocks and all successive differences}
Write $n=2m\nu+j$. Lemma~\ref{spectral} gives
\begin{equation}\label{signed}
 (-1)^\nu F_m(n)=\e^{\mathsf T}B^n u_j.
\end{equation}
For $0\le j\le2m-2$, $n$ and $j$ have the same parity, so
Lemma~\ref{positive} proves strict positivity. For $j=2m-1$,
Lemma~\ref{wave} gives zero.

For $0\le j\le2m-3$, the two terms $n$ and $n+1$ lie in
the same block. Apply \eqref{signed} to both and use the recurrence:
\begin{align}
 (-1)^\nu\bigl(F_m(n+1)-F_m(n)\bigr)
 &=\e^{\mathsf T}B^n(Bu_{j+1}-u_j)\notag\\
 &=\e^{\mathsf T}B^n u_{j+2}.\label{difference}
\end{align}
If $j\le2m-4$, Lemma~\ref{positive} applies to $t=j+2$,
giving strict increase. If $j=2m-3$, then $u_{j+2}=u_{2m-1}=0$,
giving the equality of the final two terms. Taking $m=2^k$
proves Theorem~\ref{main}.\hfill$\square$

\section*{Tasks 2--5: proposed follow-up questions}
These are proposed directions, not additional solved problems.
\begin{enumerate}[label=\textbf{Task \arabic*:},start=2,leftmargin=*]
\item Determine whether every positive signed block is log-concave:
for $b_j=(-1)^\nu f_{m,2m\nu+j}$, is
$b_j^2\ge b_{j-1}b_{j+1}$ at every interior index? Identify all
equality cases or the first counterexample.
\item Determine exact $2$-adic valuations of the nonzero block
entries and their positive differences. Can the path representation
explain cancellations modulo powers of two and yield formulas in
the binary digits of $n$?
\item For even $m$ that are not powers of two, classify the negative
terms and failures of monotonicity of the original $f_{m,n}(-1)$.
Its sign sequence is no longer the single square wave in \eqref{F};
investigate a decomposition into finite-path contributions.
\item Obtain uniform asymptotic formulas for the walk sums in
\eqref{signed} and \eqref{difference} when $m\to\infty$,
$n/m^2$ stays in a compact subset of $(0,\infty)$, and the
position $j/(2m)$ varies, including the endpoints of a block.
\end{enumerate}

\section*{Verification and source scope}
The accompanying standard-library program \texttt{verify.py} uses only
integer arithmetic. Its four groups check 2,256 Lucas-parity cases,
5,884 central-vector identities, 3,160 independent binomial-sum/walk
identities, and 3,160 sign/zero cases together with 2,960 exact difference
identities. The proof above, rather than these finite checks, establishes
the universal quantifiers.

All 22 pages of the current preprint and all 26 pages of the published
final were read. The preprint states the target on page 18; the final
states it on page 22. Their earlier results include the trigonometric
formula, positivity for sufficiently late blocks, the first block,
and terminal equality. Those results are credited to Dilcher and Ulas.
The present argument supplies positivity and strict monotonicity in
every block by a nonnegative path representation.

The MathDB API's reference text labels a different harmonic-sum
statement as refuted. That text is not a resolution of the target,
which remains a conjecture in both current source versions. The
database's solution and attempt fields were empty at the recorded
checks. Bounded literature and public claim-index searches found no
prior full claim for this target. This is an authoring-agent proof and
audit; no independent peer review or formal verification is claimed.

\begin{thebibliography}{9}
\bibitem{DU}
K. Dilcher and M. Ulas,
\emph{Divisibility and Arithmetic Properties of a Class of Sparse Polynomials},
Integers \textbf{22} (2022), A8, 26 pp.
\href{https://math.colgate.edu/~integers/w8/w8.pdf}{Published final};
Conjecture 2, p.~22; Proposition 7, pp.~18--20.
\bibitem{DUa}
K. Dilcher and M. Ulas, same title,
\href{https://arxiv.org/abs/2008.13475}{arXiv:2008.13475v2},
24 September 2021, 22 pp.; Conjecture 5.8, p.~18.
\bibitem{DB}
\href{https://mathdb.com/p/347788}{MathDB 347788},
\emph{Positivity and monotonicity conjecture for the sequence
$f_{2^k,n}$}, checked 18 September 2026.
\end{thebibliography}
\end{document}
